\exercice

Dans un repère orthonormé, on considère les points $((nA))((A))$, $((nB))((B))$ et $((nC))((C))$.

\begin{enumerate}
  \item
\begin{enumerate}
  \item Déterminer les coordonnées de $I$, milieu de $[((nA))((nC))]$.
  \item Déterminer les coordonnées de $((nD))$, symétrique de $((nB))$ par rapport à $I$.
  \item En déduire que $((nABCD))$ est un parallélogramme.
\end{enumerate}
\item
\begin{enumerate}
  \item Calculer la longueur $((nB))((nC))$.
  \item On admet que $((nA))((nB))=((AB[0]))\sqrt{((AB[1]))}$ et $((nA))((nC))=((AC[0]))\sqrt{((AC[1]))}$. Le triangle $((nA))((nB))((nC))$ est-il équilatéral ? isocèle en $((nB))$ ? rectangle en $((nB))$ ?
\end{enumerate}
\item Préciser, si possible, la nature du parallélogramme $((nABCD))$.
\end{enumerate}
